% GENERATED FILE. Edit canonical lessons, then run npm run generate:books.
% canonical-source-hash: fnv1a64:65e103c80fd1aa73
% canonical-lessons: KA-C541-kasarattu, KA-C541-sahisnute, KA-C541-parvatarohana, KA-C541-billugarike, KA-C541-huridumbisu

\chapter{Fitness}
\label{ch:ka-fitness}
\hlchaptermodality{541}

\begin{chapteropening}
\textbf{By the end of this chapter:} \emph{I can say a workout, endurance, mountaineering, archery and to cheer on.}

Complete the last lesson of chapter 541: I can say a workout, endurance, mountaineering, archery and to cheer on.
\end{chapteropening}

\section[\texorpdfstring{\kn{ಕಸರತ್ತು} (kasarattu)}{ಕಸರತ್ತು (kasarattu)}]{\kn{ಕಸರತ್ತು} (kasarattu) — a workout; physical exercise}

\label{lesson:KA-C541-kasarattu}



\begin{quote}
Before the new one: say the Kannada for to kick, then the Kannada for to punch.
\end{quote}

\subsection*{You'll want to know: \kn{ಕಸರತ್ತು}}
\textbf{\kn{ಕಸರತ್ತು}} — \emph{kasarattu} — \textquotedblleft{}a workout; physical exercise\textquotedblright{}.

\begin{quote}
\kn{ಪ್ರತಿದಿನ} \kn{ಬೆಳಗ್ಗೆ} \kn{ಅಪ್ಪ} \kn{ಅರ್ಧ} \kn{ಗಂಟೆ} \kn{ಕಸರತ್ತು} \kn{ಮಾಡುತ್ತಾರೆ}.
\end{quote}

\begin{quote}
\emph{pratidina beḷagge appa ardha gaṇṭe kasarattu māḍuttāre.}
\end{quote}

\begin{quote}
Every morning father works out for half an hour.
\end{quote}

Also a clever trick or stunt.

\begin{cousinweb}
\textbf{\kn{ಕಸರತ್ತು}} (\emph{kasarattu}) is usually traced to Persian \emph{kasrat}, \textquotedblleft{}exercise\textquotedblright{}, which reached Kannada through Hindustani.
\end{cousinweb}

\begin{grammarlens}[title={What you've built}]
One more word for describing what you lived and felt.
\end{grammarlens}

\subsection*{Your turn}
\begin{itemize}
\raggedright
  \item \emph{Say it:} \emph{kasarattu}
  \item \emph{Say it:} \emph{kasarattu}, once more
  \item \emph{Say it:} \emph{pratidina beḷagge appa ardha gaṇṭe kasarattu māḍuttāre.} — Every morning father works out for half an hour.
  \item \emph{Recall:} say the Kannada for to kick, then the Kannada for to punch, then say \emph{kasarattu} again
\end{itemize}

\begin{tcolorbox}[breakable,colback=teal!4,colframe=teal!35!black,title={Before you move on}]
What does \kn{ಕಸರತ್ತು} mean? (\textquotedblleft{}A workout; physical exercise\textquotedblright{}.) Say it once more.
\end{tcolorbox}
\section[\texorpdfstring{\kn{ಸಹಿಷ್ಣುತೆ} (sahiṣṇute)}{ಸಹಿಷ್ಣುತೆ (sahiṣṇute)}]{\kn{ಸಹಿಷ್ಣುತೆ} (sahiṣṇute) — endurance; tolerance}

\label{lesson:KA-C541-sahisnute}



\begin{quote}
Before the new one: say the Kannada for to punch, then the Kannada for a workout.
\end{quote}

\subsection*{You'll want to know: \kn{ಸಹಿಷ್ಣುತೆ}}
\textbf{\kn{ಸಹಿಷ್ಣುತೆ}} — \emph{sahiṣṇute} — \textquotedblleft{}endurance; tolerance\textquotedblright{}.

\begin{quote}
\kn{ದೂರ} \kn{ಓಡಲು} \kn{ಸಹಿಷ್ಣುತೆ} \kn{ಬೇಕು}.
\end{quote}

\begin{quote}
\emph{dūra ōḍalu sahiṣṇute bēku.}
\end{quote}

\begin{quote}
Running far takes endurance.
\end{quote}

Also tolerance of other views.

\begin{cousinweb}
\textbf{\kn{ಸಹಿಷ್ಣುತೆ}} (\emph{sahiṣṇute}) is borrowed from Sanskrit \emph{sahiṣṇutā}, from \emph{sahiṣṇu}, \textquotedblleft{}bearing, enduring\textquotedblright{} (Monier-Williams).
\end{cousinweb}

\begin{grammarlens}[title={What you've built}]
One more word for describing what you lived and felt.
\end{grammarlens}

\subsection*{Your turn}
\begin{itemize}
\raggedright
  \item \emph{Say it:} \emph{sahiṣṇute}
  \item \emph{Say it:} \emph{sahiṣṇute}, once more
  \item \emph{Say it:} \emph{dūra ōḍalu sahiṣṇute bēku.} — Running far takes endurance.
  \item \emph{Recall:} say the Kannada for to punch, then the Kannada for a workout, then say \emph{sahiṣṇute} again
\end{itemize}

\begin{tcolorbox}[breakable,colback=teal!4,colframe=teal!35!black,title={Before you move on}]
What does \kn{ಸಹಿಷ್ಣುತೆ} mean? (\textquotedblleft{}Endurance; tolerance\textquotedblright{}.) Say it once more.
\end{tcolorbox}
\section[\texorpdfstring{\kn{ಪರ್ವತಾರೋಹಣ} (parvatārōhaṇa)}{ಪರ್ವತಾರೋಹಣ (parvatārōhaṇa)}]{\kn{ಪರ್ವತಾರೋಹಣ} (parvatārōhaṇa) — mountaineering}

\label{lesson:KA-C541-parvatarohana}



\begin{quote}
Before the new one: say the Kannada for a workout, then the Kannada for endurance.
\end{quote}

\subsection*{You'll want to know: \kn{ಪರ್ವತಾರೋಹಣ}}
\textbf{\kn{ಪರ್ವತಾರೋಹಣ}} — \emph{parvatārōhaṇa} — \textquotedblleft{}mountaineering\textquotedblright{}.

\begin{quote}
\kn{ನನ್ನ} \kn{ಅಣ್ಣನಿಗೆ} \kn{ಪರ್ವತಾರೋಹಣ} \kn{ತುಂಬಾ} \kn{ಇಷ್ಟ}.
\end{quote}

\begin{quote}
\emph{nanna aṇṇanige parvatārōhaṇa tumbā iṣṭa.}
\end{quote}

\begin{quote}
My elder brother loves mountaineering.
\end{quote}

\kn{ಪರ್ವತ}, mountain + \kn{ಆರೋಹಣ}, climbing.

\begin{cousinweb}
\textbf{\kn{ಪರ್ವತಾರೋಹಣ}} (\emph{parvatārōhaṇa}) is a learned compound built from Sanskrit \emph{parvata}, \textquotedblleft{}rock, mountain\textquotedblright{} (CDIAL 7945), and \emph{ārohaṇa}, \textquotedblleft{}ascending\textquotedblright{} (Monier-Williams).
\end{cousinweb}

\begin{grammarlens}[title={What you've built}]
One more word for describing what you lived and felt.
\end{grammarlens}

\subsection*{Your turn}
\begin{itemize}
\raggedright
  \item \emph{Say it:} \emph{parvatārōhaṇa}
  \item \emph{Say it:} \emph{parvatārōhaṇa}, once more
  \item \emph{Say it:} \emph{nanna aṇṇanige parvatārōhaṇa tumbā iṣṭa.} — My elder brother loves mountaineering.
  \item \emph{Recall:} say the Kannada for a workout, then the Kannada for endurance, then say \emph{parvatārōhaṇa} again
\end{itemize}

\begin{tcolorbox}[breakable,colback=teal!4,colframe=teal!35!black,title={Before you move on}]
What does \kn{ಪರ್ವತಾರೋಹಣ} mean? (\textquotedblleft{}Mountaineering\textquotedblright{}.) Say it once more.
\end{tcolorbox}
\section[\texorpdfstring{\kn{ಬಿಲ್ಲುಗಾರಿಕೆ} (billugārike)}{ಬಿಲ್ಲುಗಾರಿಕೆ (billugārike)}]{\kn{ಬಿಲ್ಲುಗಾರಿಕೆ} (billugārike) — archery}

\label{lesson:KA-C541-billugarike}



\begin{quote}
Before the new one: say the Kannada for endurance, then the Kannada for mountaineering.
\end{quote}

\subsection*{You'll want to know: \kn{ಬಿಲ್ಲುಗಾರಿಕೆ}}
\textbf{\kn{ಬಿಲ್ಲುಗಾರಿಕೆ}} — \emph{billugārike} — \textquotedblleft{}archery\textquotedblright{}.

\begin{quote}
\kn{ಅವಳು} \kn{ಬಿಲ್ಲುಗಾರಿಕೆಯಲ್ಲಿ} \kn{ಪದಕ} \kn{ಗೆದ್ದಳು}.
\end{quote}

\begin{quote}
\emph{avaḷu billugārikeyalli padaka geddaḷu.}
\end{quote}

\begin{quote}
She won a medal in archery.
\end{quote}

\kn{ಬಿಲ್ಲು}, a bow + -\kn{ಗಾರಿಕೆ}, the craft of.

\begin{cousinweb}
\textbf{\kn{ಬಿಲ್ಲುಗಾರಿಕೆ}} (\emph{billugārike}) is a Kannada compound of \kn{ಬಿಲ್ಲು}, \textquotedblleft{}a bow\textquotedblright{}, and -\emph{gārike}, \textquotedblleft{}the craft of\textquotedblright{}; the origin of the suffix is not settled.
\end{cousinweb}

\begin{grammarlens}[title={What you've built}]
One more word for describing what you lived and felt.
\end{grammarlens}

\subsection*{Your turn}
\begin{itemize}
\raggedright
  \item \emph{Say it:} \emph{billugārike}
  \item \emph{Say it:} \emph{billugārike}, once more
  \item \emph{Say it:} \emph{avaḷu billugārikeyalli padaka geddaḷu.} — She won a medal in archery.
  \item \emph{Recall:} say the Kannada for endurance, then the Kannada for mountaineering, then say \emph{billugārike} again
\end{itemize}

\begin{tcolorbox}[breakable,colback=teal!4,colframe=teal!35!black,title={Before you move on}]
What does \kn{ಬಿಲ್ಲುಗಾರಿಕೆ} mean? (\textquotedblleft{}Archery\textquotedblright{}.) Say it once more.
\end{tcolorbox}
\section[\texorpdfstring{\kn{ಹುರಿದುಂಬಿಸು} (huridumbisu)}{ಹುರಿದುಂಬಿಸು (huridumbisu)}]{\kn{ಹುರಿದುಂಬಿಸು} (huridumbisu) — to cheer on; to encourage}

\label{lesson:KA-C541-huridumbisu}



\begin{quote}
Before the new one: say the Kannada for mountaineering, then the Kannada for archery.
\end{quote}

\subsection*{You'll want to know: \kn{ಹುರಿದುಂಬಿಸು}}
\textbf{\kn{ಹುರಿದುಂಬಿಸು}} — \emph{huridumbisu} — \textquotedblleft{}to cheer on; to encourage\textquotedblright{}.

\begin{quote}
\kn{ಪ್ರೇಕ್ಷಕರು} \kn{ಆಟಗಾರರನ್ನು} \kn{ಜೋರಾಗಿ} \kn{ಹುರಿದುಂಬಿಸಿದರು}.
\end{quote}

\begin{quote}
\emph{prēkṣakaru āṭagārarannu jōrāgi huridumbisidaru.}
\end{quote}

\begin{quote}
The spectators cheered the players on loudly.
\end{quote}

Past (I …): \kn{ಹುರಿದುಂಬಿಸಿದೆ}, \emph{huridumbiside}.

The ones cheered take -\kn{ಅನ್ನು}.

\begin{cousinweb}
\textbf{\kn{ಹುರಿದುಂಬಿಸು}} (\emph{huridumbisu}) is a native Kannada verb, \textquotedblleft{}to fire up, to cheer on\textquotedblright{}; its origin is not settled.
\end{cousinweb}

\begin{grammarlens}[title={What you've built}]
One more word for describing what you lived and felt.
\end{grammarlens}

\subsection*{Your turn}
\begin{itemize}
\raggedright
  \item \emph{Say it:} \emph{huridumbisu}
  \item \emph{Say it:} \emph{huridumbisu}, once more
  \item \emph{Say it:} \emph{prēkṣakaru āṭagārarannu jōrāgi huridumbisidaru.} — The spectators cheered the players on loudly.
  \item \emph{Recall:} say the Kannada for mountaineering, then the Kannada for archery, then say \emph{huridumbisu} again
\end{itemize}

\begin{tcolorbox}[breakable,colback=teal!4,colframe=teal!35!black,title={Before you move on}]
What does \kn{ಹುರಿದುಂಬಿಸು} mean? (\textquotedblleft{}To cheer on; to encourage\textquotedblright{}.) Say it once more.
\end{tcolorbox}
